\section{实验评估}
\label{sec:evaluation}

本文从模型保真度、重排批次和重排时延三个方面比较GA-LK与现有编译方法，并通过受控对照分析多层前瞻与遗传搜索的作用。

\begin{table*}[!t]
\caption{ZAC18与QMAP154上三种编译方法的总体结果}
\label{tab:main-results}
\centering
\PaperTableSetup
\begin{PaperRevisionBlock}
\begin{tabular*}{\textwidth}{@{\extracolsep{\fill}}llrrrr@{}}
\toprule
数据集 & 方法 & 模型保真度 & 重排批次 & \shortstack{重排时延\\(ms)} & \shortstack{完整结果\\/$N$}\\
\midrule
\multirow{3}{*}{\shortstack{ZAC18\\($N_{\rm ind}=\DefaultZACStrictN$)}}
& ZAC & \pgfmathprintnumber[fixed,precision=4,zerofill]{\DefaultZACMOneF} & \DefaultZACMOneB & \DefaultZACMOneT & \DefaultZACMOneV\\
& ICCAD/QMAP & \pgfmathprintnumber[fixed,precision=4,zerofill]{\DefaultZACMTwoF} & \DefaultZACMTwoB & \DefaultZACMTwoT & \DefaultZACMTwoV\\
& GA-LK & \textbf{\boldmath\pgfmathprintnumber[fixed,precision=4,zerofill]{\DefaultZACMFourF}} & \textbf{\DefaultZACMFourB} & \textbf{\DefaultZACMFourT} & \DefaultZACMFourV\\
\midrule
\multirow{3}{*}{\shortstack{QMAP154\\($N_{\rm ind}=\DefaultQMAPStrictN$)}}
& ZAC & \pgfmathprintnumber[fixed,precision=6,zerofill]{\DefaultQMAPMOneF} & \DefaultQMAPMOneB & \DefaultQMAPMOneT & \DefaultQMAPMOneV\\
& ICCAD/QMAP & \pgfmathprintnumber[fixed,precision=6,zerofill]{\DefaultQMAPMTwoF} & \DefaultQMAPMTwoB & \DefaultQMAPMTwoT & \DefaultQMAPMTwoV\\
& GA-LK & \textbf{\boldmath\pgfmathprintnumber[fixed,precision=6,zerofill]{\DefaultQMAPMFourF}} & \textbf{\DefaultQMAPMFourB} & \textbf{\DefaultQMAPMFourT} & \DefaultQMAPMFourV\\
\bottomrule
\end{tabular*}
\end{PaperRevisionBlock}
\PaperTableNote[\textwidth]{注：保真度取几何均值，重排批次和时延取算术均值；粗体标示同一数据集内的最优值。$N_{\rm ind}$为三种方法共同有效的不同电路数，QMAP154的\PaperRevision{\DefaultQMAPStrictFileN} 个共同有效文件去重后对应\PaperRevision{\DefaultQMAPStrictN} 个电路。完整结果数按全部输入文件统计。}
\end{table*}


\subsection{实验设置}

实验采用ZAC论文中的18个电路（ZAC18）\cite{lin2025zac}和MQT QMAP示例集中的154个输入文件（QMAP154）\cite{wille2023qmap}。QMAP154经电路规范化与去重后包含151个不同电路；等价输入合并为同一分析单元，其保真度在对数域取均值。外部基线为Lin等人的ZAC\cite{lin2025zac}和Stade等人的路由感知放置ICCAD/QMAP\cite{stade2025routingaware}，后者基于MQT QMAP v3.2.0实现。三种编译器使用相同的规范化输入与门序列。

硬件模型包含$100\times100$个存储区SLM阱位、$7\times20$个纠缠区门位对和一个$100\times100$的AOD阵列。保真度采用式~\eqref{eq:fidelity}，参数为$f_1=0.9997$、$f_2=0.995$、$f_{\rm exc}=0.9975$、$f_{\rm tran}=0.999$及$T_2=1.5\times10^6\,\mu$s，与ZAC的评估设置一致\cite{lin2025zac}。每次SLM--AOD阱转移计入$N_{\rm tran}$，输运时间计入$t_q$。外部基线使用各自的原生输出，GA-LK采用第~\ref{sec:method}节的候选可行性检查。单次运行限时600~s。保真度、重排批次和时延仅统计三种编译器均完成编译且满足$t_q<T_2$的不同电路；完整结果的数量按全部输入文件统计。

程序以C++17 Release模式编译，关闭OpenMP与快速数学优化。单线程计时平台配备14核Apple M4 Pro和24~GB内存，使用macOS~26.6.2、Python~3.10.19和AppleClang~21.0.0。

主比较采用表~\ref{tab:solver-bounds}的求解预算及$(H_{\max},\alpha,\rho,\epsilon)=(8,0.5,0.7,0.05)$。GA-LK以随机种子0、1和2分别运行，逐电路取中位数。受控组件比较的参数与样本范围在对应小节给出。

\subsection{编译质量}

GA-LK在两个基准集上均提高了模型保真度，平均重排批次与重排时延也低于两种基线，三种方法的总体结果见表~\ref{tab:main-results}。在ZAC18上，保真度几何均值相对ZAC和ICCAD/QMAP分别提高\PaperPercentGain{\ZACMFourF}{\ZACMOneF}和\PaperPercentGain{\ZACMFourF}{\ZACMTwoF}；在QMAP154上，分别提高\PaperPercentGain{\QMAPMFourF}{\QMAPMOneF}和\PaperPercentGain{\QMAPMFourF}{\QMAPMTwoF}，对应增幅见图~\ref{fig:experimental-summary}(a)。

个别电路获得了更大的保真度改善。在表~\ref{tab:representative-circuits}的四个实例中，量子傅里叶变换电路\texttt{QFT-18}的相对增幅最大，三个Ising电路也均优于两种基线，逐电路增幅见图~\ref{fig:experimental-summary}(b)。这些实例来自两个基准集，各取两项；要求较高基线保真度不低于0.05，并按GA-LK相对该值的对数保真度差排序。表中列出电路规模和绝对保真度，图中给出相对各基线的增幅。

% Auto-generated by prepare_experimental_summary.py; do not edit.
\def\FigSixLabelA{\texttt{QFT-18}}
\def\FigSixLabelB{\texttt{Ising-42}}
\def\FigSixLabelC{\texttt{Ising-16}}
\def\FigSixLabelD{\texttt{Ising-13}}
\def\RepresentativeCircuitRows{%
ZAC18 & \texttt{\detokenize{qft_n18}} & 18/306 & 0.0683 & 0.0640 & \textbf{0.0894} \\%
ZAC18 & \texttt{\detokenize{ising_n42}} & 42/82 & 0.3723 & 0.3853 & \textbf{0.4167} \\%
\addlinespace[1pt]
QMAP154 & \texttt{\detokenize{ising_model_16}} & 16/150 & 0.2519 & 0.2469 & \textbf{0.2788} \\%
QMAP154 & \texttt{\detokenize{ising_model_13}} & 13/120 & 0.3432 & 0.3411 & \textbf{0.3759} \\%
}

% Display aliases only: the frozen rows, circuit identities, and numbers are unchanged.
\ExplSyntaxOn
\NewExpandableDocumentCommand{\PaperCircuitLabel}{m}{
  \str_case_e:nnF {#1} {
    {qft_n18}{QFT-18}
    {ising_n42}{Ising-42}
    {ising_model_16}{Ising-16}
    {ising_model_13}{Ising-13}
  } {#1}
}
\ExplSyntaxOff
\begin{table}[!t]
\caption{典型高增益电路的模型保真度比较}
\label{tab:representative-circuits}
\centering
\PaperTableSetup
\let\PaperRawTexttt\texttt
\renewcommand{\texttt}[1]{\PaperRawTexttt{\PaperCircuitLabel{#1}}}
\begin{tabularx}{\columnwidth}{@{}lLcrrr@{}}
\toprule
基准集 & 电路 & $n/g_2$ & ZAC & ICCAD & GA-LK\\
\midrule
\RepresentativeCircuitRows
\bottomrule
\end{tabularx}
\PaperTableNote{注：电路简称与图~\ref{fig:experimental-summary}一致；$n/g_2$为量子比特数/双比特门数。ICCAD指ICCAD/QMAP，粗体为该电路的最高模型保真度。选例条件见正文，总体结果见表~\ref{tab:main-results}。}
\end{table}

% Declare both single-column tables before the wide figure so that they remain
% beside the results discussion rather than entering the final-page float queue.
\begin{table}[!t]
\caption{前瞻层数上限与编译质量}
\label{tab:sensitivity}
\centering
\PaperTableSetup
\begin{tabular*}{\columnwidth}{@{}l@{\extracolsep{\fill}}rr@{}}
\toprule
前瞻层数上限 $H$ & 保真度比 & 重排批次比\\
\midrule
0 & \HorizonExtZeroFidelityRatio & \HorizonExtZeroBatchRatio\\
1 & \HorizonExtOneFidelityRatio & \HorizonExtOneBatchRatio\\
2 & \HorizonExtTwoFidelityRatio & \HorizonExtTwoBatchRatio\\
4 & \HorizonExtFourFidelityRatio & \HorizonExtFourBatchRatio\\
8 & \HorizonExtEightFidelityRatio & \HorizonExtEightBatchRatio\\
\bottomrule
\end{tabular*}
\PaperTableNote{注：\HorizonExtTargetN 个预定电路中，\HorizonExtCompleteN 个在各配置下均有有效模型评分，两列使用这一共同集合。每个配置取\HorizonExtSeedN 次运行的中位数，再相对$H=8$按电路计算后取几何均值。保真度比越大越好，重排批次比越小越好。}
\end{table}

\begin{figure*}[t]
  \centering
  \resizebox{\textwidth}{!}{% The same frozen aggregate and circuit values, on zero-based gain axes.
% Shared vector-figure vocabulary for the Chinese IEEE manuscript.
% Every figure in this directory inputs this file, so the manuscript only
% needs the standard TikZ package and the libraries already loaded by
% paper_zh.tex.  Keep the guard: figures may be input more than once while
% editing or when a stand-alone smoke-test document is built.
\ifdefined\GALKFigureStyleLoaded\else
\def\GALKFigureStyleLoaded{1}

% Baseline-derived visual tokens: ICCAD uses a restrained blue/orange device
% palette, while ZAC uses conventional black axes and a small fixed plot set.
\definecolor{galkInk}{HTML}{231F20}
\definecolor{galkMuted}{HTML}{666666}
\definecolor{galkGrid}{HTML}{D6D6D6}
\definecolor{galkPaper}{HTML}{FFFFFF}
% Semantic colors are fixed across every figure: orange=entanglement zone,
% blue=storage zone, black=atom identity, red=invalidity, and blue dashes=future
% prediction.  Plot and search-group colours are deliberately independent of
% the two physical-zone colours.
\definecolor{galkStorage}{HTML}{0064BC}
\definecolor{galkStorageFill}{HTML}{E8F1F7}
\definecolor{galkEntangle}{HTML}{E37323}
\definecolor{galkEntangleFill}{HTML}{F9ECE3}
\definecolor{galkResident}{HTML}{485CB3}
\definecolor{galkResidentFill}{HTML}{EEF0F8}
\definecolor{galkGood}{HTML}{A4AF07}
\definecolor{galkBad}{HTML}{C84E45}
\definecolor{galkLook}{HTML}{0064BC}
\definecolor{galkData}{HTML}{315E8A}
\definecolor{galkTail}{HTML}{B44B27}
\definecolor{galkGroupOne}{HTML}{4C78A8}
\definecolor{galkGroupTwo}{HTML}{7A7A7A}
\definecolor{galkGroupThree}{HTML}{9C755F}

\tikzset{
  % Shared engineering-figure contract.  Physical snapshots must place an
  % orange entanglement band above a blue storage field; only their scale may
  % change.  Labels are aligned at the upper-left of both regions.
  galk/base/.style={
    font=\rmfamily\fontsize{9.0pt}{10.0pt}\selectfont,
    text=galkInk,
    line cap=butt,
    line join=miter,
    >=Stealth
  },
  galk/panel title/.style={
    font=\rmfamily\bfseries\fontsize{9.2pt}{10.2pt}\selectfont,
    text=galkInk,
    anchor=west
  },
  galk/label/.style={
    font=\rmfamily\fontsize{9.0pt}{10.0pt}\selectfont,
    text=galkInk
  },
  galk/small label/.style={
    font=\rmfamily\fontsize{8.7pt}{9.7pt}\selectfont,
    text=galkInk
  },
  galk/zone label/.style={
    font=\rmfamily\bfseries\fontsize{8.7pt}{9.7pt}\selectfont,
    text=galkInk,
    anchor=north west,
    inner sep=.5pt
  },
  galk/entanglement zone/.style={
    draw=galkEntangle,
    fill=galkEntangleFill,
    line width=.40pt
  },
  galk/storage zone/.style={
    draw=galkStorage,
    fill=galkStorageFill,
    line width=.40pt
  },
  galk/site/.style={
    circle,
    draw=galkInk!68,
    fill=white,
    line width=.36pt,
    minimum size=1.35mm,
    inner sep=0pt
  },
  galk/occupied atom/.style={
    circle,
    draw=galkInk,
    fill=galkInk,
    line width=.36pt,
    minimum size=1.55mm,
    inner sep=0pt
  },
  galk/target site/.style={
    circle,
    draw=galkInk!60,
    fill=galkInk!10,
    line width=.34pt,
    minimum size=1.55mm,
    inner sep=0pt
  },
  galk/q label/.style={
    font=\rmfamily\itshape\fontsize{8.5pt}{9.5pt}\selectfont,
    text=galkInk,
    inner sep=.25pt
  },
  galk/rydberg pair/.style={
    draw=galkInk!62,
    densely dashed,
    line width=.36pt
  },
  galk/accepted path/.style={
    -{Stealth[length=.95mm,width=.62mm]},
    draw=galkInk,
    line width=.48pt
  },
  galk/alternative path/.style={
    -{Stealth[length=.95mm,width=.62mm]},
    draw=galkMuted,
    densely dashed,
    line width=.48pt
  },
  galk/future path/.style={
    -{Stealth[length=.95mm,width=.62mm]},
    draw=galkLook,
    densely dashed,
    line width=.48pt
  },
  galk/invalid path/.style={
    -{Stealth[length=.95mm,width=.62mm]},
    draw=galkBad,
    densely dashed,
    line width=.50pt
  }
}
\fi

\providecommand{\FigSixDataRoot}{figures/data}
% Auto-generated by prepare_experimental_summary.py; do not edit.
\def\FigSixZACN{18}
\def\FigSixQMAPN{120}
\def\FigSixZACWTL{12/0/6}
\def\FigSixQMAPWTL{49/0/71}
\def\FigSixZACXMax{18.5}
\def\FigSixQMAPXMax{120.5}
\def\FigSixZACYMin{-0.109360343307}
\def\FigSixZACYMax{1.07142308229}
\def\FigSixQMAPYMin{-0.134112355854}
\def\FigSixQMAPYMax{1.1203299954}
\def\FigSixQMAPTrueMin{-0.187327670958}
\def\FigSixQMAPTrueMax{9.88913154837}
\def\FigSixQMAPLowTailN{6}
\def\FigSixQMAPHighTailN{6}
\def\FigSixMechanismZACN{18}
\def\FigSixMechanismQMAPN{122}
\def\FigSixAblationHN{39}
\def\FigSixAblationGAN{38}
\def\FigSixStageN{171}
\def\FigSixFullCompileMean{35.956}
\def\FigSixAdditiveMean{35.010}

% Auto-generated by prepare_experimental_summary.py; do not edit.
\def\FigSixLabelA{\texttt{QFT-18}}
\def\FigSixLabelB{\texttt{Ising-42}}
\def\FigSixLabelC{\texttt{Ising-16}}
\def\FigSixLabelD{\texttt{Ising-13}}
\def\RepresentativeCircuitRows{%
ZAC18 & \texttt{\detokenize{qft_n18}} & 18/306 & 0.0683 & 0.0640 & \textbf{0.0894} \\%
ZAC18 & \texttt{\detokenize{ising_n42}} & 42/82 & 0.3723 & 0.3853 & \textbf{0.4167} \\%
\addlinespace[1pt]
QMAP154 & \texttt{\detokenize{ising_model_16}} & 16/150 & 0.2519 & 0.2469 & \textbf{0.2788} \\%
QMAP154 & \texttt{\detokenize{ising_model_13}} & 13/120 & 0.3432 & 0.3411 & \textbf{0.3759} \\%
}

\edef\FigSixZACvsZAC{\fpeval{100*(\ZACMFourF/\ZACMOneF-1)}}
\edef\FigSixZACvsICCAD{\fpeval{100*(\ZACMFourF/\ZACMTwoF-1)}}
\edef\FigSixQMAPvsZAC{\fpeval{100*(\QMAPMFourF/\QMAPMOneF-1)}}
\edef\FigSixQMAPvsICCAD{\fpeval{100*(\QMAPMFourF/\QMAPMTwoF-1)}}

\pgfplotsset{
  figsix axis/.style={
    axis lines=left,
    axis line style={draw=galkInk,line width=.5pt},
    tick style={draw=galkInk,line width=.45pt},
    tick align=outside,
    tick label style={font=\rmfamily\fontsize{8.5pt}{9.5pt}\selectfont},
    label style={font=\rmfamily\fontsize{9pt}{10pt}\selectfont},
    grid=none,clip=false,scaled ticks=false,
    ymin=0,ymax=45,ytick={0,10,20,30,40},
    ylabel={Fidelity improvement (\%)},
    ybar,bar width=10pt,
    nodes near coords,
    point meta=y,
    every node near coord/.style={
      font=\rmfamily\fontsize{8pt}{9pt}\selectfont,text=galkInk,
      /pgf/number format/fixed,/pgf/number format/precision=1,
      /pgf/number format/zerofill,
      anchor=south,yshift=1pt}
  },
  figsix zac/.style={draw=galkData,fill=galkData!85,line width=.45pt},
  figsix iccad/.style={draw=galkTail,fill=galkTail!65,line width=.45pt}
}

\begin{tikzpicture}[galk/base]
  \path[use as bounding box] (-.05,-.30) rectangle (17.85,5.92);
  \begin{pgfinterruptboundingbox}
  % The legend names the baseline in each direct comparison.
  \filldraw[draw=galkData,fill=galkData!85,line width=.45pt]
    (5.05,5.63) rectangle +(0.30,.16);
  \node[anchor=west,font=\rmfamily\fontsize{9pt}{10pt}\selectfont]
    at (5.43,5.71) {relative to ZAC};
  \filldraw[draw=galkTail,fill=galkTail!65,line width=.45pt]
    (8.30,5.63) rectangle +(0.30,.16);
  \node[anchor=west,font=\rmfamily\fontsize{9pt}{10pt}\selectfont]
    at (8.68,5.71) {relative to ICCAD/QMAP};

  \node[galk/panel title] at (.82,5.17)
    {(a) Overall fidelity improvement};
  \begin{axis}[
    figsix axis,at={(.98cm,.68cm)},anchor=south west,
    width=5.90cm,height=4.00cm,scale only axis,
    xmin=.55,xmax=2.45,xtick={1,2},
    xticklabels={\shortstack{ZAC18\\$N=\FigSixZACN$},
                 \shortstack{QMAP154\\$N=\FigSixQMAPN$}}
  ]
    \addplot[figsix zac,bar shift=-7.5pt,
      every node near coord/.append style={anchor=south east,xshift=1pt}]
      coordinates {(1,\FigSixZACvsZAC) (2,\FigSixQMAPvsZAC)};
    \addplot[figsix iccad,bar shift=7.5pt,
      every node near coord/.append style={anchor=south west,xshift=-1pt}]
      coordinates {(1,\FigSixZACvsICCAD) (2,\FigSixQMAPvsICCAD)};
  \end{axis}

  \node[galk/panel title] at (8.27,5.17)
    {(b) Selected circuit improvements};
  \begin{axis}[
    figsix axis,at={(8.42cm,.68cm)},anchor=south west,
    width=8.92cm,height=4.00cm,scale only axis,
    xmin=.45,xmax=4.55,xtick={1,2,3,4},
    xticklabels={\shortstack{\FigSixLabelA\\ZAC18},
                 \shortstack{\FigSixLabelB\\ZAC18},
                 \shortstack{\FigSixLabelC\\QMAP154},
                 \shortstack{\FigSixLabelD\\QMAP154}},
    bar width=8pt
  ]
    \draw[draw=galkGrid,line width=.4pt]
      (axis cs:2.5,0) -- (axis cs:2.5,43);
    \addplot[figsix zac,bar shift=-5pt,
      every node near coord/.append style={anchor=south east,xshift=1pt}]
      table[x=x,y expr={100*(\thisrow{ratio_zac}-1)}]
      {\FigSixDataRoot/fig6_selected_cases.dat};
    \addplot[figsix iccad,bar shift=5pt,
      every node near coord/.append style={anchor=south west,xshift=-1pt}]
      table[x=x,y expr={100*(\thisrow{ratio_iccad}-1)}]
      {\FigSixDataRoot/fig6_selected_cases.dat};
  \end{axis}
  \end{pgfinterruptboundingbox}
\end{tikzpicture}
}
  \caption{GA-LK相对ZAC和ICCAD/QMAP的模型保真度改善。(a) 两个基准集上的几何均值增幅，分别使用18和120个共同有效的不同电路。(b) 表~\ref{tab:representative-circuits}所列四个高增益电路的逐项增幅。纵轴均为$100(F_{\rm GA\text{-}LK}/F_{\rm baseline}-1)$，从零起画；数值越大表示改善越多。}
  \label{fig:experimental-summary}
\end{figure*}

重排批次减少意味着原子输运所需的串行阶段更少。GA-LK在ZAC18上的平均批次数由ZAC的\ZACMOneB 和ICCAD/QMAP的\ZACMTwoB 降至\ZACMFourB；在QMAP154上，由\QMAPMOneB 和\QMAPMTwoB 降至\QMAPMFourB。平均重排时延也分别降至\ZACMFourT\,ms和\QMAPMFourT\,ms，均低于两种基线。GA-LK的平均输运开销在批次数与持续时间两个方面均得到改善。

\subsection{多层评价与遗传搜索}
\label{sec:ablation}

加入后续门层的信息后，$H_{\max}=8$的模型保真度几何均值较零视界配置提高\PaperPercentGain{\AblationHRatio}{1}。在相同目标函数与评价预算下，遗传搜索较逐坐标确定性贪心及局部精修提高\PaperPercentGain{\AblationGARatio}{1}。

视界比较采用$H_{\max}=8$与$H_{\max}=0$，共用$(\alpha,\rho,\epsilon)=(0.5,0.7,0.05)$、搜索预算和存储位候选域。两者均满足当前层的物理约束。最大视界为8层的配置估计后续门层的物理代价和窗口末端的回迁代价，并依据未来复用确定驻留保护集合；零视界配置以当前层代价为目标。遗传搜索与贪心方法的比较固定$H_{\max}=8$目标及全部物理约束。

对照使用全部ZAC18电路和30个分层选取的QMAP154电路。QMAP154按双比特门数分为不超过300、301--1500和大于1500三个层级，各取10个；层内按规范化电路内容的哈希顺序选取，清单见代码仓库。视界与搜索策略比较分别纳入\AblationHN 和\AblationGAN 个具备完整三种子结果的电路。

\subsection{参数选择与编译开销}

候选数量与前瞻深度决定了搜索需要的计算量。评价预算、存储位域、最大视界和衰减权重对保真度与编译时间的影响见表~\ref{tab:sensitivity}。比较采用固定的12个电路：ZAC18按量子比特数分为不超过32、33--64和大于64三个层级，QMAP154沿用双比特门数分层，各层按相同哈希规则取2个电路。

在所评估配置中，降低评价预算或缩小存储位域后，保真度仍接近默认值，编译时间比分别降至\SensitivityBudgetLowTimeRatio 和\SensitivityReturnSmallTimeRatio。最大视界缩短至4层时，保真度也接近默认配置。

候选搜索与未来状态评价增加了GA-LK的编译计算量。固定12电路子集上的端到端编译时间中位数为\StrictMFourTime\,s，相对ICCAD/QMAP的逐电路时间比几何均值为\StrictMFourVsMTwo。各编译器使用原生单线程实现，计时为预热后三次运行的中位数。
